% ============================================
% Supplementary Note 2: Assembly History and Origin-Correlated Species Fates
% ============================================
\subsection*{Supplementary Note 2: Assembly history and origin-correlated species fates}
\label{sec:suppl:note2}

\paragraph{Assembly history.}
To understand the mechanistic basis of community-level selection, we compared coalescence outcomes to a ``direct assembly'' control. In coalescence, two parental communities are first assembled separately from the species pool, then mixed; in direct assembly, all species from both pools are mixed simultaneously without prior assembly history.

During assembly, species compete and some go extinct. Those that survive to coexistence tend to have weaker mutual interactions on average. We quantified this assembly effect by measuring the average pairwise interaction coefficient $\langle \alpha_{ij} \rangle$ among species at different stages (Supplementary Fig.~18). Across all tested interaction-strength parameter values ($\mu = 0$ to $1.2$), post-assembly communities exhibited substantially reduced average pairwise interaction coefficients compared to the pre-assembly pool. This reduction becomes more pronounced at higher $\mu$, where stronger competition drives more species to extinction, leaving only those with weaker mutual interactions. The assembly process thus creates internal coherence within each parental community: species that survived together share compatible (less competitive) interactions.

We hypothesized that this shared history would increase Dominance probability compared to direct assembly. Simulations supported this prediction (Extended Data Fig.~7): across interaction-strength parameter values $\mu = 0.2$ to $1.0$, coalescence consistently produced higher Dominance fractions than direct assembly. At $\mu = 0.2$, coalescence yielded 37\% Dominance versus 12\% for direct assembly; at $\mu = 0.6$, 59\% versus 20\%; at $\mu = 1.0$, 65\% versus 22\%. Correspondingly, direct assembly showed higher Restructuring rates, as expected when species from different pools have not been pre-filtered for compatibility.

Experimental results corroborate these findings (Supplementary Fig.~19). Under Base and Nutr$+$ conditions, coalescence showed elevated Dominance rates compared to direct assembly, consistent with simulation predictions. The Nutr$-$ condition showed more variable results, likely due to weaker competitive interactions in nutrient-limited environments.

\paragraph{Pairwise selection correlation.}
To quantify whether species from the same parental community exhibit correlated selection during coalescence, we developed a pairwise selection correlation metric. For each coalescence event, species were assigned to their parental community of origin (\rev{the parental community with higher abundance} if present in both). For each species pair, we evaluated whether their presence/absence patterns in the offspring showed shared fate (both present or both absent) or different fate (one survives, one extinct).

The shared-fate rate was converted to a correlation metric: $\rho = 2 \times \text{shared-fate rate} - 1$, yielding values from $-1$ (all pairs have different fates) to $+1$ (all pairs share fate). We computed $\rho_{\text{same}}$ for within-community pairs and $\rho_{\text{cross}}$ for cross-community pairs, then averaged across coalescence events. To test significance, we generated null distributions by shuffling species origin labels (1,000 permutations) and computed permutation $p$-values for $\Delta = \bar{\rho}_{\text{same}} - \bar{\rho}_{\text{cross}}$.

Positive $\bar{\rho}_{\text{same}}$ indicates species from the same parental community share fates (survive or go extinct together); negative $\bar{\rho}_{\text{cross}}$ indicates cross-community species have opposite fates. In both simulations (Extended Data Fig.~5a--d) and experiments (Extended Data Fig.~6), we observed $\bar{\rho}_{\text{same}} > 0$ and $\bar{\rho}_{\text{cross}} < 0$, with $\Delta$ increasing with interaction strength. At weak interactions ($\mu = 0.3$), both within-community and cross-community pairs show positive correlations with only a small difference ($\Delta = 0.05$; Extended Data Fig.~5a). As interaction strength increases, the difference becomes pronounced: at $\mu = 0.6$, within-community pairs show positive correlation while cross-community pairs become negative ($\Delta = 0.58$; Extended Data Fig.~5b); at $\mu = 0.8$, the separation is even larger ($\Delta = 0.94$; Extended Data Fig.~5c). This supports the interpretation that assembly history can structure pair-level species fates during coalescence, providing an outcome-level hallmark of community-level selection.
